\section{局限性}
\label{sec:limitations-zh}

程序化语料合法且可复现，但只代表一种设计分布，不能覆盖当代作曲实践的广度。模型可能学习生成器规律，而无法迁移到独立的人类创作乐谱。当前尚未在合法提供的 MusicXML 上进行外部验证，也尚未收集作曲家或后调性分析者的盲评。

正式矩阵是单种子描述性实验。它没有估计不同训练初始化、语料种子或候选采样流之间的变异，也没有报告推断显著性检验。大部分约束指标同时参与候选选择，因此其 $K=4$ 改善反映已实现的选择机制，而不是独立感知标准。隐藏密度曲线是部分例外，但变化幅度较小。

自动指标把复杂音乐判断压缩为可检查特征。音级集合覆盖率类似召回率，音程向量忽略顺序与音区，循环序列准确率依赖序列化音符顺序，姿态一致性则使用人工设定的特征目标。所有神经配置的严格序列形式准确率为 0，这表明较高聚合完成率不能被理解为成功实现十二音序列。

模型还会低估请求跨度。虽然 MusicXML 的小节数与声部数完全正确，拟议模型的平均内容跨度比仅为 0.5195，因为导出阶段对末尾空时间进行了补齐。因此，结构有效性不等于材料持续覆盖全部请求时长，也不能证明出版级记谱质量或作曲工作流中的实际用途。

事件词表采用 4/4 拍、十六分音符网格，不包含连音、微分音、变拍号、扩展技法、详细演奏法、正式评估中的力度以及制谱布局。一个乐器标签会复制到各声部，因此当前接口不支持混合编制。256 标记上下文、固定四候选预算、固定惩罚权重，以及缺少权重和采样温度敏感性研究，也进一步限制了本文结论。
